\chapter{Build and start}\label{build-and-start}
\index{Building the emulator}
\index{Startup}

Requirements: a POSIX host, a C17 compiler, CMake 3.20 or newer, GNU Readline development headers/library, and a terminal. The linked Readline runtime library must also be available when running a copied executable. Neither Python nor TeX is required for an ordinary build or binary-distribution execution. Optional developer verification and media regeneration have separate dependencies.

From the repository root:

\begin{commandexample}[language=sh]{Host shell: build and start}
cmake -S . -B build -DCMAKE_BUILD_TYPE=Release
cmake --build build
cd build
cp -R ../doc/usermanual/examples .
./sim4000 -a
\end{commandexample}

CMake supplies \passthrough{\lstinline!RC4000.cfg!}, \passthrough{\lstinline!RC4000.dfc403.cfg!}, \passthrough{\lstinline!MONITOR!}, \passthrough{\lstinline!UNIT0!}, and \passthrough{\lstinline!UNIT0.original!} in the build directory and creates \passthrough{\lstinline!exchange!}. The default configuration selects the native \passthrough{\lstinline!UNIT0!} with \passthrough{\lstinline!fromhost!} and \passthrough{\lstinline!tohost!}. The unchanged historical image retains \passthrough{\lstinline!fromdos!} and \passthrough{\lstinline!todos!}; select it with \passthrough{\lstinline!./sim4000 -a --disk UNIT0.original!}.

The command \passthrough{\lstinline!-a!} starts the interactive interface, requests disk startup, and immediately runs. Normal disk writes are volatile: they are visible during this emulator session but disappear when it exits. Choose a persistence mode before starting if you want to keep your work.

The copied \texttt{examples/} directory is optional for normal operation, but
contains the input files used in the guest tutorials. It is not executed
automatically. Keep it beside the runtime files if you move your working
directory; the examples then need no source checkout.

You are now in the running emulator, not at the host shell. Continue directly
with \manualref{first-monitor-session}{First monitor session} to answer the
date question. Skip \manualref{run-from-your-own-working-directory}{the optional
working-directory setup} unless you want to run a separate copy.
